\S\ref{subsec:w3} reports that our rate sweep cannot measure the exponent
Theorem~\ref{thm:general} predicts, and withdraws two earlier readings of
it. This appendix gives the sweep, the two faults, and the registered rule
applied to the result, so that a reader can check the withdrawal rather
than take it.

\paragraph{The registered rule.} Theorem~\ref{thm:general} predicts
distortion decaying as $2^{-2R/\deff}$, a slope in a narrow band on a
log-rate axis. We registered the band as $[-2.4, -1.6]$, to hold on at
least three of four base models, and swept quantizer width $b \in
\{1,2,3,4,8,16\}$ with the fit over $b \in \{2,3,4,8\}$, $b = 1$ being
clipping-dominated. The rule voids the earlier falsification if the slope
lies in the band on three of four bases, and lets it stand if
$|\mathrm{slope}| < 0.5$ on three of four.

\paragraph{First fault: a mismatched reference.} An earlier version
reported the prediction falsified on measured slopes of $-0.10$ and
$-0.21$. \textbf{Those two numbers are withdrawn.} Each finite-rate curve
was compared against a $b = \infty$ reference computed at a
\emph{different} ridge value and aggregated over a different number of
cohorts, so the difference was not taken within one method. The symptom was
visible in the output at the time and we did not act on it: on two bases
the excess at finite rate came out \emph{below} its own asymptote, which
the model does not permit.

\begin{table}[t]
\centering
\small
\caption{Rate sweep against matched asymptotes: same ridge $\lambda$, same
realisation, same cohort for the curve and its $b = \infty$ reference, with
the equality asserted from the configuration files before any fit is
attempted. Worst-task NLL excess in nats, at four decimals because at three
the reader cannot check the two facts that matter. \textbf{Ten of the
twenty-four measured values fall below their own asymptote}, which the model
forbids; all but one are smaller than $10^{-4}$ nats and are consistent with
numerical noise, the exception being Llama at $b=3$ at $-0.0020$. The fit is
ordinary least squares of $\log_2(\text{excess} - \text{asymptote})$ on $b$
over $b \in \{2,3,4,8\}$, keeping residuals above $10^{-9}$; the ``fit''
column names the points that survived that filter, and both surviving fits
rest partly on residuals below this metric's own resolution of about
$0.001$ nats (\S\ref{sec:discussion}).}
\label{tab:w3}
\begin{tabular}{lccccccccccc}
\toprule
base & $b{=}1$ & $b{=}2$ & $b{=}3$ & $b{=}4$ & $b{=}8$ & $b{=}16$ &
$b{=}\infty$ & fit & slope & $R^2$ \\
\midrule
Llama-3.1-8B & 0.6329 & 0.0967 & 0.0803 & 0.0845 & 0.0823 & 0.0823 & 0.0823 & 2,4,8 & $-2.10$ & 0.9843 \\
Mistral-7B & 0.2106 & 0.0497 & 0.0480 & 0.0471 & 0.0474 & 0.0473 & 0.0474 & 2,3 & --- & --- \\
Qwen2.5-7B & 0.0746 & 0.0110 & 0.0102 & 0.0101 & 0.0101 & 0.0100 & 0.0102 & 2,3 & --- & --- \\
Yi-1.5-9B & 0.2099 & 0.0413 & 0.0361 & 0.0357 & 0.0356 & 0.0355 & 0.0356 & 2,3,4 & $-3.50$ & 0.9996 \\
\bottomrule
\end{tabular}
\end{table}

\paragraph{Second fault: neither surviving fit is supported at the
precision the metric provides,} and the published three-decimal table
concealed it. The fits are on the residual above the asymptote, and those
residuals are mostly tiny. Llama's three fitted points have residuals
$0.0144$, $0.0023$ and $0.0000029$ nats. The last is the $b = 8$ point and
is what makes the fit possible at all: without it Llama has two usable
points and is unfittable like Mistral and Qwen. It sits about $350$ times
below this metric's stated resolution of $0.001$ nats, so on a $\log_2$
axis it carries enormous leverage while being, by our own account of the
metric, indistinguishable from zero. Yi is worse in one respect: two of its
three fitted residuals, $0.00046$ and $0.000044$, are below the resolution.

We also withdraw this section's earlier account of why Mistral and Qwen are
unfittable. It said each ``reaches its asymptotic value inside the fit
window'', which is benign and wrong. What happens is that their $b=4$ and
$b=8$ points fall \emph{below} their asymptotes and are removed by the
residual filter: the same anomaly described above, which the published
caption asserted no longer occurred. It occurs on all four bases, ten times
in twenty-four measurements, and Table~\ref{tab:w3} now reports it.

\paragraph{The rule, applied as written.} The registered rule needs the
slope in $[-2.4, -1.6]$ on three of four bases to void the falsification,
and $|\mathrm{slope}| < 0.5$ on three of four for it to stand. We obtain
one and zero, so the outcome is \textsc{unresolved}. We apply the rule to
the fits as computed rather than adjusting the residual filter now that we
can see which points it admits: moving a threshold after seeing its
consequences is what the pre-registration exists to prevent, and here it
would move the verdict in our own favour.

The outcome is unresolved for a more basic reason than the rule
contemplates, though. The rule presumes fits worth judging, and on this
grid there are none: only $b \in \{1,2\}$ is comfortably above the noise
floor on any base, and $b=1$ is excluded as clipping-dominated. The
registered window $\{2,3,4,8\}$ is too coarse for bases that reach their
noise floor by three bits. A finer low-rate grid would settle it, and we
have deliberately not chosen one after seeing which grid would produce a
fit.

\paragraph{What survives is operational.} With the ridge on, three bits
buys essentially the whole effect: from $b = 3$ upward every base is within
$0.005$ nats of its own asymptote, and sixteen bits adds nothing measurable
over four. Only $b = 1$ is far off. A practitioner quantizing a merged
adapter can stop at three or four bits, and that recommendation does not
depend on any exponent.

\paragraph{Caveats.} This sweep is a single initialisation draw on the
shared-initialisation cohort, so it says nothing about whether finite-rate
behaviour differs by regime; repeating it on an independent cohort asks a
different question from the one registered here and would need its own
rule. Both faults were ours and both were in the analysis rather than the
cells: the first in the reference the analyser selected, the second in the
precision at which we reported and read the result. The second was
invisible at three decimals and would have stayed invisible, which is why
the table above prints the fitted points and the residual scale rather than
rounded values alone.
